%% kapa: exam=1 round=34 year=2023 period=1 form=A subject=real_estate subject_ko=부동산학원론 seq=24 number=104
%% source: Q-Net 감정평가사 기출문제 (cst00309) artlSeq=5211789 file=raw/1008/5211789/2023년도 제34회 감정평가사 1차시험 1교시.pdf
%% conversion: kordoc 4.24.51 -> tools/gen1.py (auto); review=proofread-sonnet; flags=none
\begin{question}{104}
부동산시장세분화에 관한 설명으로 옳지 않은 것은?
\begin{choices}
  \item 시장세분화는 가격차별화, 최적의사결정, 상품차별화 등에 기초하여 부동산시장을 서로 다른 둘 또는 그 이상의 상위시장으로 묶는 과정이다.
  \item 시장을 세분화하는데 주로 사용되는 기준으로는 지리적 변수, 인구통계학적 변수, 심리적 변수, 행동적 변수 등이 있다.
  \item 시장세분화전략은 세분된 시장을 대상으로 상품의 판매 지향점을 명확히 하는 것을 말한다.
  \item 부동산회사가 세분시장을 평가할 때, 우선해야 할 사항으로 적절한 시장규모와 성장성을 들 수 있다.
  \item 세분시장에서 경쟁력과 매력도를 평가할 때 기존 경쟁자의 위협, 새로운 경쟁자의 위협, 대체재의 위협, 구매자의 협상력 증가 위협, 공급자의 협상력 증가 위협 등을 고려한다.
\end{choices}
\end{question}
\answer{1}
